\documentclass[12pt, a4paper]{article}
\usepackage[utf8]{inputenc}
\usepackage[T1]{fontenc}
\usepackage{amsmath, amssymb}
\usepackage{geometry}
\usepackage{booktabs}
\usepackage{titlesec}
\usepackage{xcolor}
\usepackage{mdframed}
\usepackage{enumitem}
\usepackage{array}
\usepackage{parskip}
\usepackage{listings}
\usepackage{courier}

\geometry{margin=2.5cm, top=3cm, bottom=3cm}

\definecolor{ballblue}{RGB}{0, 82, 147}
\definecolor{lightgray}{RGB}{245, 245, 245}
\definecolor{darkgray}{RGB}{80, 80, 80}
\definecolor{accentgray}{RGB}{180, 180, 180}
\definecolor{codegreen}{RGB}{40, 120, 40}
\definecolor{codegray}{RGB}{100, 100, 100}
\definecolor{codebg}{RGB}{250, 250, 250}
\definecolor{codeblue}{RGB}{0, 60, 180}
\definecolor{codeorange}{RGB}{180, 80, 0}

\titleformat{\section}
  {\large\bfseries\color{ballblue}}{}
  {0em}{}[\color{ballblue}\titlerule]
\titlespacing{\section}{0pt}{18pt}{8pt}

\titleformat{\subsection}
  {\normalsize\bfseries\color{darkgray}}{}
  {0em}{}
\titlespacing{\subsection}{0pt}{12pt}{4pt}

\mdfdefinestyle{destaque}{
  backgroundcolor=lightgray,
  linecolor=ballblue, linewidth=1.2pt,
  leftline=true, rightline=false, topline=false, bottomline=false,
  innerleftmargin=10pt, innerrightmargin=10pt,
  innertopmargin=8pt, innerbottommargin=8pt,
  skipabove=6pt, skipbelow=6pt,
  nobreak=true,
}

\mdfdefinestyle{codigo}{
  backgroundcolor=codebg,
  linecolor=accentgray, linewidth=0.8pt,
  leftline=true, rightline=false, topline=false, bottomline=false,
  linecolor=codegreen,
  innerleftmargin=10pt, innerrightmargin=10pt,
  innertopmargin=8pt, innerbottommargin=8pt,
  skipabove=6pt, skipbelow=6pt,
  nobreak=true,
}

\lstdefinelanguage{Julia}{
  keywords={using, function, for, in, if, else, end, return, true, false, @variable, @constraint, @objective, Model, optimize!, Min},
  keywordstyle=\color{codeblue}\bfseries,
  comment=[l]{\#},
  commentstyle=\color{codegray}\itshape,
  stringstyle=\color{codeorange},
  morestring=[b]",
  basicstyle=\ttfamily\scriptsize,
  breaklines=true,
  showstringspaces=false,
}

\lstset{language=Julia}

% =============================================================================
\begin{document}

% ── Capa ─────────────────────────────────────────────────────────────────────
\begin{center}
  {\color{ballblue}\rule{\linewidth}{2pt}}\\[8pt]
  {\LARGE\bfseries Modelo de Programa\c{c}\~{a}o Linear Estoc\'{a}stica}\\[4pt]
  {\large\color{darkgray} Planejamento Macrosc\'{o}pico --- Ball Corporation}\\[8pt]
  {\color{ballblue}\rule{\linewidth}{2pt}}
\end{center}

\vspace{10pt}
\begin{center}
  \small\color{darkgray}
  \textbf{1.~Minimiza\c{c}\~{a}o de custos}
  \enspace\textbullet\enspace
  \textbf{2.~Conjuntos}
  \enspace\textbullet\enspace
  \textbf{3.~Par\^{a}metros}
  \enspace\textbullet\enspace
  \textbf{4.~Vari\'{a}veis}
  \enspace\textbullet\enspace
  \textbf{5.~A F\'{i}sica do Modelo}
  \enspace\textbullet\enspace
  \textbf{6.~Fun\c{c}\~{a}o Objetivo}
  \enspace\textbullet\enspace
  \textbf{7.~Restri\c{c}\~{o}es}
  \enspace\textbullet\enspace
  \textbf{8.~Formulac\c{c}\~{a}o Compacta}
  \enspace\textbullet\enspace
  \textbf{9.~Implementa\c{c}\~{a}o}
\end{center}
\vspace{4pt}
{\color{accentgray}\rule{\linewidth}{0.5pt}}

% ============================================================
\section{Modelo de minimiza\c{c}\~{a}o de custos}

O modelo busca \textbf{minimizar o custo total de opera\c{c}\~{a}o} da malha industrial da Ball
Corporation. Al\'{e}m de contemplar os custos de produ\c{c}\~{a}o, estocagem e transfer\^{e}ncia,
o modelo incorpora penalidades de risco: a multa por quebra de N\'{i}vel de Servi\c{c}o (venda perdida)
e a penalidade por opera\c{c}\~{a}o abaixo do Estoque de Seguran\c{c}a. O modelo \'{e} estritamente cont\'{i}nuo
e cobre um horizonte de 15 dias.

\[
  \boxed{
    \begin{aligned}
      \min \; Z \;=\;\;
      & \underbrace{\sum_{p,l,i,t} k_{i,p,l} \cdot x_{i,p,l,t}}_{\text{produ\c{c}\~{a}o}}
        \;+\;
        \underbrace{\sum_{p,l,i,t} v_{i,p,l} \cdot e_{i,p,l,t}}_{\text{estoque}} 
        \;+\;
        \underbrace{\sum_{i,o,d,t} ct_{o,d} \cdot \dfrac{f_{i,o,d,t}}{200{.}000}}_{\text{log\'{i}stica}}
      \\[12pt]
      & +\;
        \underbrace{\sum_{p,l,i,t} M \cdot w_{i,p,l,t}}_{\text{venda perdida}}
        \;+\;
        \underbrace{\sum_{p,l,i,t} c_{ss} \cdot \text{falta\_ss}_{i,p,l,t}}_{\text{risco de estoque}}
    \end{aligned}
  }
\]

\begin{center}
  \small\color{darkgray}
  Sujeito \`{a}s restri\c{c}\~{o}es f\'{i}sicas e operacionais R1--R10 detalhadas na Se\c{c}\~{a}o 7.
\end{center}

% ============================================================
\section{Conjuntos (\'Indices)}

\begin{center}
\renewcommand{\arraystretch}{1.4}
\begin{tabular}{>{\bfseries\color{ballblue}}c p{8cm} c r}
  \toprule
  S\'{i}mbolo & Descri\c{c}\~{a}o & \'{I}ndice & Cardinalidade \\
  \midrule
  $I$           & SKUs (produtos finais) & $i$ & 30 \\
  $P$           & Plantas (f\'{a}bricas) & $p,\,o,\,d$ & 10 \\
  $L_p$         & Linhas da planta $p$ & $l$ & 3 por planta \\
  $T$           & Per\'{i}odos $(t_1, \ldots, t_{15})$ & $t$ & 15 dias \\
  $\mathcal{C}$ & Triplas cap\'{a}veis $(i,p,l)$: formato de $i$ = formato de $l \in L_p$ & $(i,p,l)$ & $\subseteq I \times L$ \\
  \bottomrule
\end{tabular}
\end{center}

% ============================================================
\section{Par\^{a}metros e Fatores de Risco}

As bases estruturais (custos, demandas $\eta_{i,p,l,t}$, limites $e^{\max}_{p,l}$ e taxas de produ\c{c}\~{a}o de placa $r_{p,l,t}$) mant\^{e}m-se conforme a extra\c{c}\~{a}o dos arquivos \texttt{CSV}. Para aproximar a matem\'{a}tica do ch\~{a}o de f\'{a}brica, introduzimos:

\begin{align*}
  \text{OEE}  &: \text{Fator de Efici\^{e}ncia Global do Equipamento (Assumido: 85\%).} \\
  DM_{i,p,l}  &: \text{Demanda m\'{e}dia di\'{a}ria do SKU } i \text{ na linha } l \text{ durante o horizonte de 15 dias.} \\
  Dias_{ss}   &: \text{Meta de cobertura para Estoque de Seguran\c{c}a (Assumido: 1 dia de } DM \text{).}
\end{align*}

% ============================================================
\section{Vari\'{a}veis de Decis\~{a}o}

S\~{a}o oito vari\'{a}veis de decis\~{a}o, todas cont\'{i}nuas e n\~{a}o-negativas:

\vspace{6pt}
\begin{center}
\renewcommand{\arraystretch}{1.6}
\begin{tabular}{>{\bfseries\color{ballblue}}c p{9.5cm} l}
  \toprule
  Vari\'{a}vel & Descri\c{c}\~{a}o & Dom\'{i}nio \\
  \midrule
  $x_{i,p,l,t}$ & Volume \textbf{produzido} do SKU $i$ na linha $l \in L_p$, no dia $t$. & $\mathbb{R}_+$ \\[4pt]
  $e_{i,p,l,t}$ & N\'{i}vel de \textbf{estoque} do SKU $i$ na linha $l \in L_p$ ao final do dia $t$. & $\mathbb{R}_+$ \\[4pt]
  $a_{i,p,l,t}$ & Volume de \textbf{atendimento local} do pedido do SKU $i$ no dia $t$. & $\mathbb{R}_+$ \\[4pt]
  $f_{i,o,d,t}$ & Volume \textbf{transferido log\'{i}sticamente} da planta $o$ para a planta $d$. & $\mathbb{R}_+$ \\[4pt]
  $g_{i,p,l,t}$ & Volume \textbf{recebido internamente} pela empilhadeira da doca para o armazém. & $\mathbb{R}_+$ \\[4pt]
  $s_{i,p,l,t}$ & Volume \textbf{expedido internamente} do armazém para a doca de carregamento. & $\mathbb{R}_+$ \\[4pt]
  $w_{i,p,l,t}$ & Volume de \textbf{venda perdida} (demanda n\~{a}o atendida). & $\mathbb{R}_+$ \\[4pt]
  $\text{falta\_ss}_{i,p,l,t}$ & Volume \textbf{abaixo da meta} de Estoque de Seguran\c{c}a. & $\mathbb{R}_+$ \\
  \bottomrule
\end{tabular}
\end{center}

% ============================================================
\section{A F\'{i}sica do Modelo: Modelando o Mundo Real}

Modelos de Programa\c{c}\~{a}o Linear puros sofrem do que chamamos de \textit{A Fal\'{a}cia da M\'{e}dia}. O algoritmo otimizador tende a ser ingenuamente otimista: ele enxerga o tempo como blocos perfeitos e as m\'{a}quinas como rochas inquebr\'{a}veis. Para garantir validade t\'{a}tica, inserimos duas restri\c{c}\~{o}es avan\c{c}adas:

\vspace{6pt}
\begin{mdframed}[style=destaque]
\textbf{\color{ballblue}1. A Ilus\~{a}o da M\'{a}quina Perfeita (O Fator OEE na restri\c{c}\~{a}o R4)}

\medskip
Se a taxa manual de uma m\'{a}quina \'{e} de 1.000 latas/hora, um solver cl\'{a}ssico assumir\'{a} que entregar\'{a} 24.000 latas amanh\~{a}. Na vida real, a f\'{a}brica enfrenta a troca de turno dos operadores, limpezas, trocas de bobina e falhas el\'{e}tricas. Ao multiplicar a taxa m\'{a}xima por um $\text{OEE} = 0.85$ na restri\c{c}\~{a}o R4, for\c{c}amos a matem\'{a}tica a prever que a m\'{a}quina perder\'{a} 15\% do dia. Se n\~{a}o ensin\'{a}ssemos isso ao solver, ele venderia uma capacidade fantasma, colidindo com a Teoria das Filas no ch\~{a}o de f\'{a}brica.
\end{mdframed}

\vspace{8pt}
\begin{mdframed}[style=destaque]
\textbf{\color{ballblue}2. O Medo do Imprevisto (Estoque de Seguran\c{c}a como Soft Constraint)}

\medskip
A equa\c{c}\~{a}o f\'{i}sica R2 decreta que a f\'{a}brica inicie com os galp\~{o}es inteiramente vazios ($e^0 = 0$). Queremos que o algoritmo atue como um gerente de f\'{a}brica prudente, que constr\'{o}i um Estoque de Seguran\c{c}a ao longo do tempo para se blindar da vari\^{a}ncia de pedidos.

No entanto, se impusermos uma \textbf{Restri\c{c}\~{a}o R\'{i}gida} (Hard Constraint) como $e \ge DM$, o modelo ir\'{a} colapsar (\textit{Infeasible}) no dia $t=1$, pois \'{e} f\'{i}sicamente imposs\'{i}vel nascer com latas no estoque inicial zero.

\medskip
\textbf{A solu\c{c}\~{a}o (Restri\c{c}\~{a}o R10):} Injetamos a vari\'{a}vel de risco $\text{falta\_ss}$. 
\[ e_{i,p,l,t} + \text{falta\_ss}_{i,p,l,t} \geq \text{Meta de Seguran\c{c}a} \]
Isso d\'{a} empatia ao modelo: se ele n\~{a}o tiver estoque, a equa\c{c}\~{a}o n\~{a}o quebra, mas a vari\'{a}vel $\text{falta\_ss}$ assume o risco e paga uma pequena multa.
\end{mdframed}

% ============================================================
\section{Fun\c{c}\~{a}o Objetivo}

\[
  \begin{aligned}
    \min \; Z \;=\;\;
    & \sum_{p \in P}\sum_{\substack{l \in L_p \\ (i,p,l)\,\in\,\mathcal{C}}}\sum_{i \in I}\sum_{t \in T}
        k_{i,p,l} \cdot x_{i,p,l,t}
      \;+\;
      \sum_{p \in P}\sum_{l \in L_p}\sum_{i \in I}\sum_{t \in T}
        v_{i,p,l} \cdot e_{i,p,l,t} \\[10pt]
    & +\;
      \sum_{i \in I}\sum_{o \in P}\sum_{\substack{d \in P \\ d \,\neq\, o}}\sum_{t \in T}
        ct_{o,d} \cdot \frac{f_{i,o,d,t}}{200{.}000}
      \;+\;
      \sum_{p,l,i,t} 1.000.000 \cdot w_{i,p,l,t}
      \;+\;
      \sum_{p,l,i,t} 5 \cdot \text{falta\_ss}_{i,p,l,t}
  \end{aligned}
\]

\textbf{O Comportamento do Otimizador:} Ao precificarmos o cancelamento do pedido ($w$) em R\$ 1.000.000 e a viola\c{c}\~{a}o do estoque de seguran\c{c}a ($\text{falta\_ss}$) em meros R\$ 5, ensinamos a intelig\^{e}ncia algor\'{i}tmica a ter \textbf{empatia comercial}: O cliente sempre vem primeiro, mas, havendo folga fabril, encha o armazém para n\~{a}o pagar a taxa de risco.

% ============================================================
\newpage
\section{Restri\c{c}\~{o}es F\'{i}sicas}

\begin{mdframed}[style=destaque]
\textbf{\color{ballblue}R1 --- Limite de armazenagem.}
\medskip

O espa\c{c}o f\'{i}sico do armazém \'{e} compartilhado entre todos os SKUs:
\[
  \sum_{i \in I} e_{i,p,l,t} \;\leq\; e^{\max}_{p,l}
  \qquad \forall\, p \in P,\; l \in L_p,\; t \in T
\]
\end{mdframed}

\vspace{8pt}

\begin{mdframed}[style=destaque]
\textbf{\color{ballblue}R2 --- Balan\c{c}o de estoque.}
\medskip

Conserva\c{c}\~{a}o de massa do armazém acoplada no tempo ($e_{i,p,l,0} = 0$):
\[
  e_{i,p,l,t} = e_{i,p,l,t-1} + x_{i,p,l,t} + g_{i,p,l,t} - s_{i,p,l,t} - a_{i,p,l,t}
  \qquad \forall\, i, p, l, t \in T
\]
\end{mdframed}

\vspace{8pt}

\begin{mdframed}[style=destaque]
\textbf{\color{ballblue}R3 --- Atendimento \`{a} demanda.}
\medskip

O atendimento local $a$ somado \`{a} venda perdida $w$ iguala a demanda $\eta$:
\[
  a_{i,p,l,t} + w_{i,p,l,t} = \eta_{i,p,l,t}
  \qquad \forall\, i, p, l, t \in T
\]
\end{mdframed}

\vspace{8pt}

\begin{mdframed}[style=destaque]
\textbf{\color{ballblue}R4 --- Capacidade produtiva (Ajustada com OEE).}
\medskip

A soma das horas consumidas por todos os SKUs na linha $l$ sofre a a\c{c}\~{a}o da inefici\^{e}ncia mec\^{a}nica global antes de bater no teto do turno de 24 horas:
\[
  \sum_{i \in I} \frac{x_{i,p,l,t}}{r_{p,l,t} \cdot \text{OEE}} \;\leq\; 24
  \qquad \forall\, p \in P,\; l \in L_p,\; t \in T
\]
\end{mdframed}

\vspace{8pt}

\begin{mdframed}[style=destaque]
\textbf{\color{ballblue}R5 --- Capacidade log\'{i}stica de expedi\c{c}\~{a}o.}
\medskip

Cada planta expede no m\'{a}ximo 80 caminh\~{o}es por dia (sa\'{i}da):
\[
  \sum_{i \in I} \sum_{\substack{d \in P \\ d \neq p}}
  \frac{f_{i,p,d,t}}{200{.}000} \;\leq\; 80
  \qquad \forall\, p \in P,\; t \in T
\]
\end{mdframed}

\vspace{8pt}

\begin{mdframed}[style=destaque]
\textbf{\color{ballblue}R6 e R7 --- N\'{o} de transbordo da f\'{a}brica (docas).}
\medskip

O frete n\~{a}o entra m\'{a}gicamente no estoque; ele \'{e} balizado pelas docas ($g$ e $s$):
\[
  \sum_{l \in L_p} g_{i,p,l,t} = \sum_{\substack{o \in P \\ o \neq p}} f_{i,o,p,t} \tag{R6}
\]
\[
  \sum_{l \in L_p} s_{i,p,l,t} = \sum_{\substack{d \in P \\ d \neq p}} f_{i,p,d,t} \tag{R7}
\]
\end{mdframed}

\vspace{8pt}

\begin{mdframed}[style=destaque]
\textbf{\color{ballblue}R8 --- Capabilidade de formato e R9 --- N\~{a}o-negatividade.}
\[
  x_{i,p,l,t} = 0 \qquad \forall\, (i,p,l) \notin \mathcal{C},\; t \in T \tag{R8}
\]
Todas as 8 matrizes de vari\'{a}veis adotam $Var \geq 0$ (R9).
\end{mdframed}

\vspace{8pt}

\begin{mdframed}[style=destaque]
\textbf{\color{ballblue}R10 --- Soft Constraint de Estoque de Seguran\c{c}a.}
\medskip

O estoque de final de dia mais a complac\^{e}ncia do risco s\~{a}o cobrados:
\[
  e_{i,p,l,t} + \text{falta\_ss}_{i,p,l,t} \geq DM_{i,p,l} \cdot Dias_{ss} \qquad \forall\, (i,p,l) \in \mathcal{C},\; t \in T
\]
\end{mdframed}

% ============================================================
\newpage
\section{Formula\c{c}\~{a}o Compacta do PL}

\begin{mdframed}[backgroundcolor=blue!5, linecolor=ballblue, linewidth=2pt,
  innertopmargin=14pt, innerbottommargin=14pt,
  innerleftmargin=14pt, innerrightmargin=14pt]

\textbf{\color{ballblue}min}
\vspace{-5pt}
\[
  \begin{aligned}
    Z \;=\;\;
    & \sum_{p,l,i,t} k_{i,p,l} \cdot x_{i,p,l,t}
      \;+\; \sum_{p,l,i,t} v_{i,p,l} \cdot e_{i,p,l,t}
      \;+\; \sum_{i,o,d,t} ct_{o,d} \cdot \tfrac{f_{i,o,d,t}}{200{.}000} \\
      & +\; \sum_{p,l,i,t} M \cdot w_{i,p,l,t}
      \;+\; \sum_{p,l,i,t} c_{ss} \cdot \text{falta\_ss}_{i,p,l,t}
  \end{aligned}
\]

\textbf{\color{ballblue}s.a.}
\vspace{-15pt}
\begin{alignat}{3}
  &\textstyle\sum_{i} e_{i,p,l,t}
    &&\leq e^{\max}_{p,l}
    &&\quad\forall\,p,l,t \tag{R1}\\[4pt]
  &e_{i,p,l,t} = e_{i,p,l,t-1} + x_{i,p,l,t} + g_{i,p,l,t} - s_{i,p,l,t} - a_{i,p,l,t}
    &&
    &&\quad\forall\,i,p,l,t \tag{R2}\\[4pt]
  &a_{i,p,l,t} + w_{i,p,l,t}
    &&= \eta_{i,p,l,t}
    &&\quad\forall\,i,p,l,t \tag{R3}\\[4pt]
  &\textstyle\sum_{i} \tfrac{x_{i,p,l,t}}{r_{p,l,t}\cdot \text{OEE}}
    &&\leq 24
    &&\quad\forall\,p,l,t \tag{R4}\\[4pt]
  &\textstyle\sum_{i,d\neq p} \tfrac{f_{i,p,d,t}}{200{.}000}
    &&\leq 80
    &&\quad\forall\,p,t \tag{R5}\\[4pt]
  &\textstyle\sum_{l \in L_p} g_{i,p,l,t}
    &&= \textstyle\sum_{o \neq p} f_{i,o,p,t}
    &&\quad\forall\,i,p,t \tag{R6}\\[4pt]
  &\textstyle\sum_{l \in L_p} s_{i,p,l,t}
    &&= \textstyle\sum_{d \neq p} f_{i,p,d,t}
    &&\quad\forall\,i,p,t \tag{R7}\\[4pt]
  &x_{i,p,l,t} = 0
    &&
    &&\quad\forall\,(i,p,l)\notin\mathcal{C},\,t \tag{R8}\\[4pt]
  &e_{i,p,l,t} + \text{falta\_ss}_{i,p,l,t}
    &&\geq DM_{i,p,l} \cdot Dias_{ss}
    &&\quad\forall\,(i,p,l)\in\mathcal{C},\,t \tag{R10}\\[4pt]
  &x,e,a,f,g,s,w,\text{falta\_ss} \geq 0
    &&
    &&\quad\forall\,i,p,l,o,d,t \tag{R9}
\end{alignat}
\end{mdframed}

\vspace{8pt}
{\color{accentgray}\rule{\linewidth}{0.5pt}}
\begin{center}
  \small\color{darkgray}
  Aproximadamente $\mathbf{139{.}500}$ vari\'{a}veis cont\'{i}nuas no total.
\end{center}

% ============================================================
\newpage
\section{Implementa\c{c}\~{a}o em Julia/JuMP}

O c\'{o}digo abaixo mapeia a l\'{o}gica estoc\'{a}stica. Note o uso de \texttt{mapa\_tempo} para blindar o solver contra saltos cronol\'{o}gicos de final de m\^{e}s, garantindo o balan\c{c}o perfeito de R2.

\begin{mdframed}[style=codigo]
\begin{lstlisting}
using JuMP, Gurobi, CSV, DataFrames, Printf, Dates

# ================= PAINEL =================
fator_oee      = 0.85  # <--- Meta de Eficiencia (Gargalo f\'isico)
dias_seguranca = 1.0   # <--- Meta de Cobertura de Estoque

const GRB_ENV = Gurobi.Env()

function resolver_cenario(cenario, df_log)
    caminho = "C:\\Dados\\ocupacao_" * cenario * "\\"
    df_need = CSV.read(caminho * "input_production_need_artificial.csv", DataFrame)
    df_rate = CSV.read(caminho * "input_production_rate_artificial.csv", DataFrame)
    df_costs = CSV.read(caminho * "input_sku_costs_artificial.csv", DataFrame)
    df_limit = CSV.read(caminho * "input_line_storage_limit_artificial.csv", DataFrame)

    # ---> MAPEAMENTO CRONOLOGICO (Previne erros em viradas de mes)
    datas_demanda = [split(string(d), " ")[1] for d in df_need.deadline if !ismissing(d)]
    datas_taxa    = [split(string(d), " ")[1] for d in df_rate.ref_date if !ismissing(d)]
    datas_unicas  = sort(unique(vcat(datas_demanda, datas_taxa)))
    mapa_tempo    = Dict(data => t for (t, data) in enumerate(datas_unicas))
    
    function get_t(data_str)
        d_limpa = split(string(data_str), " ")[1]
        return get(mapa_tempo, d_limpa, -1)
    end

    df_need.prod_need = parse_br_float.(df_need.prod_need)

    I = unique(df_need.product_id)
    P = unique(df_rate.plant)
    T = 1:length(datas_unicas)
    L_p = Dict(p => unique(df_rate[df_rate.plant .== p, :production_line]) for p in P)

    k = Dict((r.product_id, r.plant, r.production_line) => parse_br_float(r.production_cost) for r in eachrow(df_costs))
    v = Dict((r.product_id, r.plant, r.production_line) => parse_br_float(r.inventory_cost) for r in eachrow(df_costs))
    ct = Dict((r.origem, r.destino) => parse_br_float(r.custo_total_frete) for r in eachrow(df_log))
    e_max = Dict((r.plant, r.production_line) => parse_br_float(r.storage_limit) for r in eachrow(df_limit))

    eta = Dict((i,p,l,t) => 0.0 for i in I for p in P for l in L_p[p] for t in T)
    for r in eachrow(df_need)
        t = get_t(r.deadline)
        chave = (r.product_id, r.plant, r.production_line, t)
        haskey(eta, chave) && (eta[chave] += r.prod_need)
    end

    r_taxa = Dict((p,l,t) => 0.0 for p in P for l in L_p[p] for t in T)
    for r in eachrow(df_rate)
        t = get_t(r.ref_date)
        r_taxa[(r.plant, r.production_line, t)] = parse_br_float(r.rate)
    end
        
    format_linha = Dict(r.production_line => strip(lowercase(string(r.sku_size_shape))) for r in eachrow(df_rate))
    format_sku   = Dict(r.product_id => strip(lowercase(string(r.sku_size_shape))) for r in eachrow(df_need))

    C = Set((i, p, l) for i in I, p in P, l in L_p[p] 
        if get(format_sku, i, "none") == get(format_linha, l, "none_L"))

    # DM (Demanda Media Diaria)
    demanda_media = Dict((i,p,l) => sum(get(eta, (i,p,l,t), 0.0) for t in T) / 15.0 for i in I for p in P for l in L_p[p])

    modelo = Model(() -> Gurobi.Optimizer(GRB_ENV))
    set_silent(modelo)

    @variable(modelo, x[i in I, p in P, l in L_p[p], t in T] >= 0)
    @variable(modelo, e[i in I, p in P, l in L_p[p], t in T] >= 0)
    @variable(modelo, a[i in I, p in P, l in L_p[p], t in T] >= 0)
    @variable(modelo, f[i in I, o in P, d in P, t in T; o != d] >= 0)
    @variable(modelo, g[i in I, p in P, l in L_p[p], t in T] >= 0)
    @variable(modelo, s[i in I, p in P, l in L_p[p], t in T] >= 0)
    @variable(modelo, w[i in I, p in P, l in L_p[p], t in T] >= 0)
    @variable(modelo, falta_ss[i in I, p in P, l in L_p[p], t in T] >= 0) # <--- NOVO (R10)

    @objective(modelo, Min,
        sum(get(k,(i,p,l),99.0) * x[i,p,l,t] for p in P for l in L_p[p] for i in I for t in T if (i,p,l) in C) +
        sum(get(v,(i,p,l),99.0) * e[i,p,l,t] for p in P for l in L_p[p] for i in I for t in T) +
        sum(get(ct,(o,d),99999.0) * (f[i,o,d,t] / 200_000) for i in I for o in P for d in P for t in T if o != d) +
        sum(1_000_000 * w[i,p,l,t] for p in P for l in L_p[p] for i in I for t in T) +
        sum(5.0 * falta_ss[i,p,l,t] for p in P for l in L_p[p] for i in I for t in T)
    )

    @constraint(modelo, R1[p in P, l in L_p[p], t in T],
        sum(e[i,p,l,t] for i in I) <= get(e_max, (p,l), 0.0))

    @constraint(modelo, R2[i in I, p in P, l in L_p[p], t in T],
        e[i,p,l,t] == (t == 1 ? 0.0 : e[i,p,l,t-1]) + x[i,p,l,t] + g[i,p,l,t] - s[i,p,l,t] - a[i,p,l,t])

    @constraint(modelo, R3[i in I, p in P, l in L_p[p], t in T],
        a[i,p,l,t] + w[i,p,l,t] == get(eta, (i,p,l,t), 0.0))

    # ---> R4 ajustada pelo OEE
    @constraint(modelo, R4[p in P, l in L_p[p], t in T],
        sum(x[i,p,l,t] / (get(r_taxa, (p,l,t), 0.001) * fator_oee)
        for i in I if (i,p,l) in C && get(r_taxa, (p,l,t), 0.0) > 0) <= 24.0)

    @constraint(modelo, R5[p in P, t in T],
        sum(f[i,p,d,t] / 200_000 for i in I, d in P if d != p) <= 80)

    @constraint(modelo, R6[i in I, p in P, t in T],
        sum(g[i,p,l,t] for l in L_p[p]) == sum(f[i,o,p,t] for o in P if o != p))

    @constraint(modelo, R7[i in I, p in P, t in T],
        sum(s[i,p,l,t] for l in L_p[p]) == sum(f[i,p,d,t] for d in P if d != p))

    @constraint(modelo, R8[i in I, p in P, l in L_p[p], t in T; (i,p,l) not-in C], 
        x[i,p,l,t] == 0)

    # ---> R10 Soft Constraint de Seguranca
    @constraint(modelo, R10[i in I, p in P, l in L_p[p], t in T; (i,p,l) in C],
        e[i,p,l,t] + falta_ss[i,p,l,t] >= demanda_media[(i,p,l)] * dias_seguranca)

    optimize!(modelo)
    
    # [...] (Processamento e Destruicao de Memoria GC.gc())
end
\end{lstlisting}
\end{mdframed}

\end{document}